\section{The publish protocol}
\label{sec:protocol}

The guarantees are joint properties of the computation and the protocol that
writes the state the next computation reads. \spec{19} says so explicitly: they
do not survive an implementation that gets this section wrong. Three
obligations are used in the proofs.

\subsection{Tag after, immediately, and always}

\begin{definition}[tagging, \spec{19.1}]
\label{def:tagging}
For each released package: the tag is created at \Head{}; \emph{after} that
package's publish succeeds, never before; \emph{immediately} after it, not
batched to the end of the run; and for every released package whatever its
publish target, including targets that upload no artefact at all.
\end{definition}

Each clause excludes a specific failure. Tagging before publishing makes the
version unrecoverable, because \cref{def:window} will treat it as released.
Batching tags to the end of a run makes a mid-run failure republish everything
already published, which most registries reject outright. Tagging only what
uploads an artefact is the subject of \cref{prop:tagall}.

\subsection{Order}

\begin{definition}[publish order, \spec{19.2}]
\label{def:order}
Let $G$ be the workspace graph restricted to \code{publish.orderKinds} over
\emph{all} packages. Publication proceeds in a topological order of $G$, ties
broken by byte-wise least name, filtered to the plan at the end.
\end{definition}

\begin{proposition}[the order exists, is total, and is deterministic]
\label{prop:order}
\Cref{def:order} is well defined and yields the same sequence on every machine
and in every conforming implementation.
\end{proposition}

\begin{proof}
$G$ is acyclic by \cref{def:graph}, so a topological order exists. Ties are
broken by a total order on names (byte-wise comparison of distinct names), so
the order is total and independent of any container's iteration order. The
filter is applied to a sequence, preserving it.
\end{proof}

\begin{remark}[order over the whole graph, then filter]
\label{rem:filter-last}
Inducing the graph on the planned packages \emph{first} is a defect, not an
optimisation. With $A \to B \to C$ where $B$ has no bump this run (the
common case in any incremental release), the induced subgraph on $\{A, C\}$
has no edge at all, so $A$ may publish against a $C$ that has not published
yet. The constraint is real even though $B$ is not released, because $A$ still
resolves $C$ through $B$ at install time. This is why the model orders over
\Pkgs{} and filters last, and why an implementation detecting a violation must
raise \diag{E197}.
\end{remark}

Order matters for a reason no later run can repair. A dependent published
before its dependency carries a range reconciled against the dependency's old
version, and both land at the same commit, so no ancestry test can tell
afterwards that anything is wrong (\spec{13.7b}). The inconsistency is
undetectable after the fact, which is why it is prevented rather than
diagnosed.

\subsection{Partial failure}

\begin{definition}[blocking, \spec{19.3}]
\label{def:blocking}
On a failure, every package in the plan that transitively depends on the failed
one over \code{publish.blockingKinds} is \emph{blocked} (\diag{W194}): planned,
but not attempted in this run. The closure is computed over the full workspace
graph, so it traverses packages that are not in the plan. Packages that are not
blocked continue to publish.
\end{definition}

Ordering and blocking use different edge-kind sets, and the asymmetry is
deliberate: \code{orderKinds} includes optional dependencies, because an
optional dependency should still be available before its dependent, while
\code{blockingKinds} excludes them, because a missing optional dependency does
not break an install and a missing required one does.

\begin{example}[an infrastructure leg]
\label{ex:terraform}
Nothing in \cref{def:tagging,def:run} requires a publish leg to be a registry
upload. A \code{terraform apply} may be such a leg: its build stage imports the
explicitly configured live resources and writes a plan, and its publish stage
applies that plan before the release tag is written. The tag records that this
version of the configuration completed successfully. It is not Terraform
resource state, does not prove that every live resource was imported, and does
not make applying an arbitrary configuration externally idempotent. Application
builds may overlap the apply; only their publication waits for the
infrastructure publish edge.
\end{example}

\begin{remark}[recorded-success abstraction]
\Cref{def:run} models confirmed successful legs by their tags. It proves
properties of that release ledger, not convergence of an external registry or
infrastructure state. A publish may take effect before its tag is recorded;
then the repository alone is insufficient and \cref{def:reconcile} or a
publish-target-specific idempotence contract is required. Terraform likewise
needs its own state or an explicit, complete import boundary before a plan can
safely describe existing resources.
\end{remark}

\begin{definition}[reconciliation, \spec{19.4}]
\label{def:reconcile}
When a publish fails because the version already exists, the implementation
verifies by registry digest or provenance attestation that the published
artefact is the one this run would have published; if it matches, the publish
counts as successful and the tag is written (\diag{W196}); if it does not
match, or identity cannot be established, the run stops with \diag{E198}.
\end{definition}

This is the one place where the model's ``a publish either succeeds or fails''
abstraction meets a third outcome, that it already happened, and
\cref{def:reconcile} maps it back onto success. Verification is what makes that
safe: without it, adopting a version published by something other than this run
would enrol a foreign artefact into the release lineage, which is a
supply-chain vulnerability rather than a convenience.
